\documentclass[10pt,a4paper]{amsart}
\usepackage[margin=19mm]{geometry}
\usepackage{amsmath,amssymb,mathtools}
\usepackage[scr=rsfs,cal=euler]{mathalfa}
\usepackage{xfrac}
\usepackage[unicode,hidelinks,bookmarks=false]{hyperref}
\newcommand{\ie}{i.e.\ }
\newcommand{\cf}{cf.\ }
\newcommand{\resp}{respectively\ }
\newcommand{\Subsection}{\subsection}
\providecommand{\nobreakdash}{\nobreak\hbox{-}}
\setlength{\emergencystretch}{2em}
\multlinegap0pt
\usepackage{xcolor}
\usepackage{amssymb}
\usepackage{float}
\usepackage{csquotes}
\usepackage{mathtools}
\newtheorem{lem}{Lemma}
\newcommand{\R}{\mathbb{R}}
\DeclareMathOperator{\SO}{SO}
\newcommand{\con}[1]{T_{#1}}
\newcommand{\frk}[1]{\mathfrak{#1}}
\DeclareMathOperator{\identity}{id}
\DeclareMathOperator{\mg}{\Omega}
\DeclareMathOperator{\nh}{\nabla_{\mathbb{H}}}
\DeclareMathOperator{\nhs}{\nabla_{\mathbb{H}^{\mg}}}
\DeclareMathOperator{\nv}{\nabla_{\mathbb{V}}}
\DeclareMathOperator{\oo}{\mathbf{\Omega}}
\title{Tensor tomography and frame flow ergodicity for magnetic flows in higher dimensions}
\author{Louis-Brahim Beaufort}
\date{Source: arXiv:2604.12495v1 (CC BY 4.0)}
\begin{document}
\maketitle

\noindent\emph{Source and licence.} Tensor tomography and frame flow ergodicity for magnetic flows in higher dimensions. Version arXiv:2604.12495v1 (CC BY 4.0), distributed by arXiv under the Creative Commons Attribution 4.0 International licence, http://creativecommons.org/licenses/by/4.0/, as recorded in the arXiv licence metadata for this exact version and shown on the abstract page https://arxiv.org/abs/2604.12495v1. Full licence notice: \texttt{licenses/arxiv-2604.12495-CC-BY-4.0.txt}.\par\smallskip

\noindent\textit{Editorial scope.} Counted unit: the article's lem lem:mag\_hor\_grad (source0.tex lines 1471--1480). The article's complete original proof of it is delivered with it (source0.tex lines 1482--1517, one \texttt{proof} environment of 36 lines). No mathematics is written, repaired or re-derived: the statement and the proof are the article's own lines, in the article's own notation, and no retained line is substituted or re-flowed.  Retained uncounted as the source-local context the proof needs: lines 768--777; lines 1264--1266; lines 1333--1337.  Nothing was omitted from any retained span, so the package carries no omission record.  No retained line contains a citation, so the wrapper carries no bibliography and no bibliography entry.  No retained line contains a figure environment, an image-inclusion command or drawing code, so no image is delivered and the package declares no asset dependency.  Editorial formatting only, in addition: an AMS \texttt{amsart} preamble that restates the article's own declarations and macros used by the retained lines, namely the environments \texttt{lem} on separate counters, together with the standard packages those lines need.  The retained environments print in this document as Lemma 1 in order of appearance, whereas the article prints the same items with the numbering of its own full document.  The complete native source of the article as distributed in the arXiv e-print for this exact version is bundled unchanged as \texttt{sources/198404/source0.tex}.
\begin{lem} \label{lem:contorsion}
For $x \in \R^n$, the contorsion tensor $\con{x}$ is given by
\begin{align*}
\con{x} = - \langle x, e_1 \rangle \oo^0 + \frac{1}{2}(\oo e_1 \wedge x - e_1 \wedge \oo x).
\end{align*}
In the decomposition $\frk{so}(n) = \frk{so}(n-1) \oplus e_1 \wedge \R^{n-1}$, this rewrites 
\begin{align} \label{eq:decomp_contorsion}
\con{x} = - \left [\langle x, e_1 \rangle \oo^0 + \frac{1}{2} x^\perp \wedge \oo e_1 \right ] - \frac{1}{2} e_1 \wedge (\oo x + \langle x, e_1 \rangle \oo e_1).
\end{align}
\end{lem}

\begin{lem} \label{lem:skew_adj}
The operators $Y_{\oo}, Y_{e_1 \wedge \oo e_1}, Y_{\oo^0}$ and $B_x^{\mg}, x \in \R^n$ are formally skew-adjoint for the standard $L^2$ inner product.
\end{lem}

\begin{align} \label{eq:skew_adjoint_part}
\sum_{j=2}^n [Y_{e_1 \wedge e_j}, Y_{T_{e_j}}]
&= n Y_{e_1 \wedge \oo e_1} - Y_{2 \oo} \\ 
&= (n - 2) Y_{e_1 \wedge \oo e_1} - 4 Y_{\oo^0}. \notag
\end{align}

\begin{lem} \label{lem:mag_hor_grad}
Let $u \in \mathcal{C}^\infty(SM)$. We have the formula 
\begin{align*}
\nhs u = \nh u - \Omega^0 \nv u.
\end{align*}
Moreover, we have the identity 
\begin{align*}
\nv^* \Omega^0 \nv u &= - \frac{n-2}{2} (\Omega v)^\vee u.
\end{align*}
\end{lem}

\begin{proof}
By definition, for $2 \leq j \leq n$, we have $B^{\mg}_{e_j} = B_{e_j} + Y_{T_{e_j}}$, where $T$ is the contorsion. By Lemma \ref{lem:contorsion}, we have the expression 
\begin{align*}
T_{e_j} = \frac{1}{2} (\oo e_1 \wedge e_j - e_1 \wedge \oo e_j).
\end{align*}

Given $u \in \mathcal{C}^\infty(SM)$, the pullback $\pi^* u$ is $\SO(n-1)$ equivariant, hence the action of the vertical field writes 
\begin{align*}
Y_{T_{e_j}} \pi^* u 
= \frac{1}{2} (0 + Y_{e_1 \wedge \oo e_j}) \pi^* u = \pi^* \frac{1}{2} \langle [\Omega w(e_j)]^\perp, \nv u \rangle = - \pi^* \frac{1}{2} \langle w(e_j), \Omega \nv u \rangle.
\end{align*}
Here, we used that $\nv u$ is valued in $\mathcal{N}$. By summing and projecting on $SM$ the above, we finally obtain 
\begin{align*}
\nhs u = \nh u - \frac{1}{2} [\Omega \nv u]^\perp = \nh u - \Omega^0 \nv u.
\end{align*}

We now prove the identity. Note that the operator $P = - \nv^* \Omega^0 \nv$ is skew-adjoint, since $\Omega^0$ is skew-symmetric. By definition, this operator is the restriction to $\pi^* \mathcal{C}^\infty(SM)$ of $Q = \sum_{j=2}^n Y_{e_1 \wedge e_j}^* Y_{T_{e_j}}$. 

By Equation \eqref{eq:skew_adjoint_part} and Lemma \ref{lem:skew_adj}, we have 
\begin{align*}
(Q - Q^*) \pi^* u &= - \sum_{j=2}^n [Y_{e_1 \wedge e_j}, Y_{T_{e_j}}] \\ 
&= (n-2) Y_{e_1 \wedge \oo e_1} \pi^* u - 4 Y_{\oo^0} \pi^* u = (n-2) \pi^* (\Omega v)^\vee u.
\end{align*}
Denote $\pi_*: \mathcal{C}^\infty(FM) \to \mathcal{C}^\infty(SM)$ the adjoint of the pullback $\pi^*: \mathcal{C}^\infty(SM) \to \mathcal{C}^\infty(FM)$; this is the fiberwise integration operator. Since $\pi^* P = Q \pi^*$ and using that $\pi_* \pi^* = \identity$ (since the fibers of $FM \to SM$ have volume $1$), we have 
\begin{align*}
\pi_* Q^* \pi^* = (Q \pi^*)^* \pi^* = P^* \pi_* \pi^* = P^*.
\end{align*}
As a result, 
\begin{align*}
2P = P - P^* = \pi_* (Q - Q^*) \pi^* = (n-2) (\Omega v)^\vee.
\end{align*}
We conclude that 
\begin{align*}
- \nv^* \Omega^0 \nv u = \frac{1}{2} P u = \frac{n-2}{2} (\Omega v)^\vee u.
\end{align*}
\end{proof}

\end{document}
